\chapter{Company Presentation: AVIONAV}

\section*{Introduction to Chapter 1}
\addcontentsline{toc}{section}{Introduction to Chapter 1}

This chapter presents the host organization where the final year project was carried out: AVIONAV. It provides an overview of the company's history, core activities, products and services, international presence, and strategic objectives. The chapter concludes by positioning the Intelligent Avionics Health Monitoring and Predictive Maintenance Platform within the context of AVIONAV's technological ecosystem and growth strategy.

\section{Company Presentation}

AVIONAV is a Tunisian company operating in the aviation and aerospace sector. Founded in 2007, the company originally established its roots in Italy before relocating its operations to Tunisia. AVIONAV is currently headquartered in Borjine, Sousse, and employs a team of approximately 10 to 50 people. Despite its relatively small size, the company distinguishes itself as one of the very few aircraft manufacturers on the African continent capable of designing and producing complete aircraft domestically.

\section{Company History}

AVIONAV was founded in 2007 with initial operations based in Italy. As part of its strategic development, the company relocated its manufacturing and design operations to Tunisia, establishing its headquarters in Borjine, Sousse. This relocation positioned AVIONAV within the Tunisian aerospace ecosystem while maintaining its international technical standards and quality requirements. Since its establishment, the company has progressively expanded its capabilities from basic aircraft assembly to full-scale aircraft design and manufacturing.

\section{Main Activities}

AVIONAV specializes in the design and manufacturing of light aircraft, typically ranging from 2 to 4 seats, as well as ultra-light and light sport aircraft (ULM/LSA). The company's production process relies on advanced manufacturing techniques including composite material fabrication, precision machining, and full aeronautical assembly. A distinctive capability of AVIONAV is its commitment to in-house manufacturing: the company produces almost every component of its aircraft internally, with engines being the only major element sourced externally. This vertical integration represents a significant achievement for a company of its size and demonstrates strong technical mastery across the aircraft production chain.

In addition to fixed-wing aircraft manufacturing, AVIONAV has expanded its scope into drone development and aeronautical systems engineering, reflecting the company's ambition to evolve with the broader aerospace industry and emerging technological trends.

\section{Products and Services}

AVIONAV's product portfolio is organized around two main aircraft families:

\begin{itemize}
    \item \textbf{Rally:} A composite aircraft constructed primarily from carbon fiber, combining lightweight construction with high structural integrity.
    \item \textbf{Storm:} An aluminum aircraft offering a more traditional construction approach, robust and suited for various operational environments.
\end{itemize}

Beyond these core aircraft families, AVIONAV has developed several noteworthy products and systems that demonstrate its engineering capabilities:

\begin{itemize}
    \item \textbf{Oxyfly:} An onboard oxygen system for aircraft, demonstrating the company's capacity to engineer integrated avionics-adjacent systems.
    \item \textbf{U Storm:} A large unmanned aerial vehicle (UAV) designed for surveillance and inspection missions, marking AVIONAV's entry into the drone market in 2023.
    \item \textbf{Helicopter Prototype:} Unveiled in 2025, representing a significant milestone in the company's engineering roadmap and expansion into rotary-wing aircraft.
\end{itemize}

\section{International Presence}

AVIONAV has established a notable international footprint, exporting dozens of aircraft to markets across Europe and the Middle East. The company's aircraft are used across a diverse range of operational contexts, including pilot training, agricultural spraying, surveillance and inspection missions, and leisure aviation. This international presence demonstrates the competitiveness of AVIONAV's products in the global light aircraft market.

In 2021, the company strengthened its international standing by acquiring Storm Aircraft, an Italian manufacturer. This strategic acquisition expanded both AVIONAV's technical know-how and its market reach. The company is also actively pursuing European airworthiness certifications to support continued export growth and compliance with international aviation standards.

\section{Strategic Objectives}

Since 2020, AVIONAV has been executing an active expansion program with the goal of producing up to 80 aircraft per year by 2026. This growth trajectory is supported by strategic plans to increase its workforce, diversify its product range, and invest in new technologies including drones and rotary-wing aircraft. The company's strategy focuses on scaling production capacity while maintaining quality standards and pursuing international certifications to access broader markets.

\section{Project Context within AVIONAV}

As AVIONAV scales its production capacity and diversifies its product portfolio, the reliability and health of its avionics systems becomes increasingly critical. The company's ambition to grow from a small domestic producer to an internationally certified aerospace manufacturer demands a corresponding elevation in its maintenance and monitoring capabilities. Traditional reactive and scheduled maintenance approaches may not be sufficient to support the planned production scale and operational complexity.

It is within this industrial and strategic context that the present project takes shape. The Intelligent Avionics Health Monitoring and Predictive Maintenance Platform developed throughout this work is intended to serve as a functional prototype for future integration into AVIONAV's technological ecosystem. By providing data-driven anomaly detection and early warning capabilities for critical avionics components, this platform aims to support AVIONAV's transition toward predictive maintenance strategies, thereby enhancing operational safety, reducing downtime, and optimizing maintenance scheduling as the company continues its growth trajectory.

\section*{Conclusion to Chapter 1}
\addcontentsline{toc}{section}{Conclusion to Chapter 1}

This chapter has presented AVIONAV, the host organization for this final year project. The company's history, core activities, product portfolio, international presence, and strategic objectives were described. The chapter positioned the Intelligent Avionics Health Monitoring and Predictive Maintenance Platform within AVIONAV's context, explaining how the prototype aligns with the company's growth strategy and its need for advanced maintenance capabilities. The next chapter will present the state of the art analysis and project positioning.